% ICLR 2026 submission.
%
% The style files in this directory are the OFFICIAL ones, downloaded from
% https://github.com/ICLR/Master-Template/tree/master/iclr2026 -- do not edit
% them. iclr2026_conference_ORIGINAL.tex is the upstream instructions file, kept
% for reference only; it is not part of this submission.
%
% Every table under tables/ is GENERATED. Do not edit them by hand:
%   python scripts/make_boba_paper_tables.py --analysis output-boba/analysis
% \todo{} marks everything not yet backed by a completed run.
%
%   cd paper && pdflatex main && bibtex main && pdflatex main && pdflatex main

\documentclass{article}
\usepackage{iclr2026_conference,times}

\usepackage{amsmath,amssymb}
\usepackage{booktabs}
\usepackage{graphicx}
\usepackage{placeins}
\usepackage{xcolor}
\usepackage{hyperref}
\usepackage{url}

\newcommand{\todo}[1]{\textcolor{red}{[TODO: #1]}}
\newcommand{\optz}{\mathrm{opt}_z}
\newcommand{\frag}{\mathrm{frag}}

\title{What Noisy Feedback Costs Bayesian\\ Optimization, and the One Number\\ That Predicts It}

% Authors must not appear in the submitted version; keep \iclrfinalcopy
% commented out until camera-ready.
\author{\todo{authors} \\
Affiliation \\
Address \\
\texttt{email}
}

%\iclrfinalcopy % Uncomment for camera-ready ONLY, never for submission.
\begin{document}

\maketitle

\begin{abstract}
Bayesian optimization driven by human feedback is usually evaluated against
regression oracles fitted to archival ratings. That entangles the cost of noisy
feedback with the cost of the oracle being wrong, and a study of one design space
cannot say what generalises. We instead inject four human-plausible error
processes into twenty analytic landscapes with published optima, standardised so
that a unit of error means the same on each: $76{,}800$ paired noisy-versus-clean
runs over four magnitudes, two onsets and twelve acquisition functions. Error as
large as the landscape's own standard deviation, present from the first rating,
costs $14.1\%$ of the available improvement; arriving at iteration 21 of 50, it
costs $2.8\%$. At that magnitude the landscape and the onset move the cost by
$8\times$ and $5\times$, and the magnitude itself by $23\times$ over its range;
the acquisition function and the kind of error move it by $2.8\times$ and
$1.3\times$. Measured rater noise
in three human studies falls where the loss is $12$--$20\%$. The expected loss
from a single noisy greedy choice, computable before any optimization is run,
explains nearly as much of the variation as the landscape descriptors and the
error magnitude together, and conditional on it the magnitude has no detectable
residual effect ($\beta = -0.07$, $p = 0.44$). Telling the surrogate the true
noise variance removes only $5\%$ of the cost, so the loss is informational;
changing the incumbent definition removes $16\%$ overall and up to $41\%$ for
one acquisition family, so robustness rankings rank acquisition-and-incumbent
pairs. When the person applies the wrong design rather than the wrong rating, a
slip of $15\%$ of each coordinate's range costs as much as one-SD rating noise,
a record of what was actually applied removes up to half of a small slip's cost,
and the design the log recommends does $1.6$--$2.6\times$ worse than the design
the person experienced. On the same grid, the fitted-oracle design reports about
a third of the cost, and improving its oracles by up to $0.29$ held-out $R^2$
does not close the gap.
\end{abstract}

\section{Introduction}

Bayesian optimization is the standard tool for optimizing an expensive
black-box function from few evaluations \citep{jones1998efficient,
shahriari2016taking, frazier2018tutorial}, and it is increasingly pointed at
objectives that only a person can evaluate: interface parameters judged for
trust or aesthetics, generated content ranked for quality, robot behaviours
rated for comfort \citep{brochu2010tutorial, gonzalez2017preferential}. The
optimizer assumes each observation is informative about the objective, but a
person is not a sensor. People rate inconsistently, drift over a session, are
systematically generous or harsh, make correlated errors across consecutive
judgements, and sometimes apply the wrong setting altogether. How much that
costs is asked routinely and, we argue, answered with a design that cannot
support the answer.

The standard construction fits a regression model to archival ratings, treats
it as the objective, injects synthetic error into its predictions, and measures
the degradation. This has two problems that no care in the optimization code
can fix. The fitted oracle is itself wrong --- in the studies motivating this
work its cross-validated $R^2$ is $0.55$ at best and negative on one dataset ---
so the measured degradation mixes the injected error with the oracle's
mis-specification, and regret is measured against a random-search estimate of
the optimum that the optimizer can legitimately beat. And the experiment runs on
one design space, so a result such as ``acquisition $X$ is the most robust'' has
no way to say whether it holds anywhere else.

We remove both problems by giving up the human. On an analytic benchmark with a
published optimum the objective is exact, regret is measured from the true
optimum, and the geometry of the landscape can be measured before the
optimization is run and
used to predict what the error will cost. Twenty such landscapes turn ``which
acquisition is robust here'' into ``which property of a problem decides whether
feedback error hurts'', a question with a transferable answer.

\paragraph{Contributions.}
\begin{enumerate}
  \item A protocol for measuring the cost of feedback error with no fitted
    oracle in the loop: twenty analytic landscapes with verified optima on a
    common scale, four human-plausible error processes, and a model-free control
    whose excess regret is exactly zero.
  \item Measurements of that cost and of what drives it
    (Sections~\ref{sec:cost}--\ref{sec:acq}). Magnitude, landscape and onset
    dominate, and the effect of magnitude is mediated by a one-shot selection
    loss computable from the landscape and the noise level in a fraction of a
    second.
  \item A matched run of the fitted-oracle design, which reports about a third
    of the cost for a reason that better oracles do not remove
    (Section~\ref{sec:fitted}).
  \item An input-error arm, in which the person applies the wrong design rather
    than giving the wrong rating: a slip of $15\%$ of the range costs as much as
    one-SD rating noise, and up to half of a small slip's cost is the record
    rather than the displacement (Section~\ref{sec:inputerror}). Six further
    checks each change one thing about the design, from the surrogate's noise
    model to the number of objectives, with a preregistered replication on
    held-out seeds (Section~\ref{sec:robustness}).
\end{enumerate}

\section{Related work}

\paragraph{Bayesian optimization under noise.}
Noise has been part of the formulation since \citet{jones1998efficient} and is
routine in practice \citep{snoek2012practical}, and several acquisitions are
built for it: augmented expected improvement \citep{huang2006global}, the
knowledge gradient \citep{frazier2009knowledge}, and noisy expected improvement
\citep{letham2019constrained}, whose Monte-Carlo form is standard in BoTorch
\citep{balandat2020botorch}. \citet{picheny2013benchmark} benchmark infill
criteria under noise on analytic problems, the closest antecedent to this work;
we differ in treating the error \emph{process} and its onset as experimental
factors, and in asking which landscape properties predict the outcome rather
than which criterion wins. Heteroscedastic surrogates \citep{kersting2007most}
address a case we do not cover: our error is homoscedastic within a run but
arrives partway through it.

\paragraph{Human-in-the-loop and preference-based optimization.}
Where the objective is a person's judgement, the usual response is to change the
observation model: elicit pairwise preferences and fit an ordinal likelihood
\citep{chu2005preference, gonzalez2017preferential}. That handles inconsistent
\emph{scaling} but not inconsistent \emph{ranking}, and it does not tell a
practitioner running a cardinal-rating study what their rater noise will cost.
Applied work that asks this directly uses the fitted-oracle simulation described
above. We trade the realism of the objective for exactness, and map measured
rater noise onto the synthetic axis afterwards.

\paragraph{Benchmarks and landscape analysis.}
Results on analytic test functions need not transfer, and suites are small and
idiosyncratic \citep{eggensperger2021hpobench}. We do not rank optimizers on a
benchmark: the twenty landscapes are twenty observations in a regression whose
unit is the landscape, and where that observational design cannot identify an
effect we manipulate the geometry directly. Predicting algorithm behaviour from
cheap landscape statistics is an old programme in evolutionary computation,
from ruggedness \citep{weinberger1990correlated} and fitness-distance correlation
\citep{jones1995fitness} to their many descendants \citep{malan2013survey}. Our
mediator belongs to that family but is defined relative to the perturbation of
interest: it is a property of the problem and the error level together, which is
what lets it predict along both axes.

\section{Setup}

\subsection{Landscapes}
We use twenty analytic benchmarks spanning $d = 2$ to $d = 11$: eight classical
functions (Schwefel, Powell, Eggholder, Ackley, Shekel, Griewank, Hartmann-3,
Hartmann-6), five further standard problems (Branin, Rosenbrock, Rastrigin,
Michalewicz, Levy), six human-performance laws including the Yerkes--Dodson
inverted-U \citep{yerkes1908relation}, and the moving-peaks benchmark
\citep{branke1999memory}. The suite was assembled for a separate study and
adopted unchanged, so its composition is not tuned to the questions asked here;
Table~\ref{tab:benchmarks} in Appendix~\ref{app:descriptors} reports its measured
geometry.

Each objective is affinely standardised to zero mean and unit variance over its
own box, using a fixed scrambled-Sobol sample of $2^{16}$ points. Raw outputs
span nearly seven orders of magnitude across the suite, so without this a single grid
of error magnitudes would be imperceptible on one landscape and overwhelming on
the next; the map is strictly increasing, so the argmax and the ordering of
designs are unchanged. One scale is deliberately left alone. Write $\optz$ for
the number of landscape standard deviations by which the optimum exceeds the
mean, which is the expected regret of a random design. It ranges from $0.76$ to
$56.8$, and that $74\times$ spread is what identifies which scale the cost of
error follows (Section~\ref{sec:currency}).

\subsection{Error processes}
\label{sec:errors}
Let $g(x)$ be the standardised objective, $t$ the iteration, $t_0$ the onset and
$\sigma_e$ the swept magnitude, in landscape standard deviations. Observations
are exact for $t \le t_0$; thereafter the optimizer sees $\tilde g_t$:
\begin{align}
  \text{\textsc{gaussian}}\quad & \tilde g_t = g(x_t) + \epsilon_t,
    & \epsilon_t &\sim \mathcal{N}(0, \sigma_e^2) \\
  \text{\textsc{bias}}\quad & \tilde g_t = g(x_t) + \sigma_e + \epsilon_t,
    & \epsilon_t &\sim \mathcal{N}(0, \sigma_e^2) \\
  \text{\textsc{drift}}\quad & \tilde g_t = g(x_t)
    + \sigma_e\tfrac{t - t_0}{T - t_0} + \epsilon_t,
    & \epsilon_t &\sim \mathcal{N}(0, \sigma_e^2) \\
  \text{\textsc{ar}(1)}\quad & \tilde g_t = g(x_t) + e_t,
    \quad e_t = \rho e_{t-1} + \eta_t,
    & \eta_t &\sim \mathcal{N}\!\left(0, \sigma_e^2(1 - \rho^2)\right)
\end{align}
with $\rho = 0.8$ and $e_{t_0+1}$ drawn from the stationary distribution, so the
AR(1) process has standard deviation $\sigma_e$ from its first step.
\textsc{bias} models a systematically generous or harsh rater; its offset is
scaled to the swept magnitude so that it and \textsc{gaussian} contrast a
systematic with a random error \emph{of the same size}, where a fixed offset
would make the two indistinguishable at the top of the grid. \textsc{drift}
models fatigue over a session, and \textsc{ar}(1) a rater whose consecutive
judgements are correlated.

Noise first affects the observation at $t_0 + 1$; the first candidate that can
react to it is proposed at $t_0 + 2$. We sweep
$\sigma_e \in \{0.05, 0.25, 1, 5\}$ and $t_0 \in \{0, 20\}$ over $T = 50$
iterations with $5$ initial quasi-random samples, twelve acquisition functions
and ten seeds: $79{,}200$ runs, of which $76{,}800$ form noisy/clean pairs.

\subsection{Optimizer and acquisitions}
Every arm uses the same surrogate, a single-task GP with a Mat\'ern kernel,
inputs normalised to the unit cube and outputs standardised, fitted by marginal
likelihood in BoTorch \citep{balandat2020botorch}. Only the acquisition function
varies. We run ten model-based arms --- expected improvement and probability of
improvement in their analytic, log \citep{ament2023unexpected} and Monte-Carlo
forms (EI, LogEI, qEI, PI, LogPI, qPI), the upper confidence bound in analytic
and Monte-Carlo form \citep{srinivas2010gaussian} (UCB, qUCB), noisy expected
improvement \citep{letham2019constrained} (qNEI), and pure exploitation of the
posterior mean (Greedy) --- and two model-free floors, uniform random sampling
and a scrambled Sobol sequence. Improvement-based acquisitions take the maximum
posterior mean at visited points as their incumbent, the choice usually
recommended under noise and, as Appendix~\ref{sec:incumbent} shows, not a
neutral one.

\subsection{Metrics}
The primary response is post-onset per-iteration excess simple regret in
landscape standard deviations: the regret of the noisy run minus that of the
identically seeded clean run, averaged over the iterations the error could have
affected. The two runs share their candidate-proposal random stream, so the
difference isolates the error rather than the seed, and because $y^\star$ is a
verified supremum rather than a random-search estimate, no regret can be
negative. For reporting we divide by $\optz$, which gives the fraction of the
achievable improvement the error destroys. We never regress that fraction on
$\optz$ or on anything collinear with it, since dividing by $\optz$ and
regressing on $\log\optz$ would manufacture the coefficient
Section~\ref{sec:currency} estimates.

\subsection{Controls}
\label{sec:controls}
\begin{itemize}
  \item \textbf{Model-free floor.} \texttt{Random} and \texttt{Sobol} never read
    an observation, so their excess regret must be identically zero, and across
    all $12{,}800$ such comparisons it is exactly zero. This negative control
    shows that no observation noise reaches the candidate stream by any route.
    The floors are then excluded from the robustness ranking, which a method
    that ignores its data would otherwise win.
  \item \textbf{Headroom.} On a landscape where the clean run already reaches
    the optimum, a null is guaranteed rather than earned. The smallest clean-run
    advantage over the model-free floor is $0.176$ (Rosenbrock), so no landscape
    is ceiling-limited.
  \item \textbf{Uniform offset.} A constant added to every observation is
    invisible to a GP that standardises its targets and uses a posterior-mean
    incumbent. Driven with a shared noise stream, the two arms propose identical designs
    to below $10^{-9}$ over the initial design and the first three model-based
    candidates (asserted by test), after which floating-point differences are
    amplified by the acquisition optimizer; in the sweep, \textsc{bias} and
    \textsc{gaussian} at onset~0 agree to within $0.012$ at every magnitude.
\end{itemize}

\section{How much does feedback error cost?}
\label{sec:cost}

Table~\ref{tab:dose} gives the cost. Error whose standard deviation equals the
landscape's own, present from the first rating, costs $14.1\%$ of the
improvement that was available, and even error five times that large costs
$22.0\%$. An optimizer that has already located a good region is hard to
dislodge: the same $1\sigma$ error arriving at iteration 21 costs $2.8\%$, a
fifth as much.

\begin{table}[h]
\centering
\caption{Fraction of the achievable improvement destroyed by feedback error,
averaged over landscapes and model-based acquisitions. Magnitudes are in
landscape standard deviations.}
\label{tab:dose}
\begin{tabular}{lrrrr}
\toprule
error present & $0.05\sigma$ & $0.25\sigma$ & $1\sigma$ & $5\sigma$ \\
\midrule
from it.\ 1 & 0.9\% & 5.3\% & 14.1\% & 22.0\% \\
from it.\ 21 & 0.3\% & 1.2\% & 2.8\% & 4.9\% \\
\bottomrule
\end{tabular}

\end{table}

Decomposing the variation by factor, at $1\sigma$ from the first rating unless
stated, puts the practical conclusions in order. Sweeping the magnitude over its
two decades moves the cost by $23\times$; changing the landscape moves it by
$8.1\times$ (Powell $0.038$ to Hartmann-6 $0.303$); moving the onset to
iteration 21, by $5.0\times$; changing the acquisition function, by $2.8\times$
(PI to qUCB); and changing the error \emph{process}, by $1.32\times$. The two
factors practitioners attend to most, which acquisition to run and which kind of
rater error to model, are the two smallest. The mean cost is $0.148$ for
\textsc{gaussian}, $0.152$ for \textsc{drift}, $0.147$ for \textsc{bias} and
$0.115$ for \textsc{ar}(1). \textsc{bias} matches \textsc{gaussian} by
construction at this onset (Section~\ref{sec:controls}) and costs $10$--$20\%$
more at the late onset, where the offset arrives as a step; correlated error is
if anything gentler, consistent with a correlated sequence being partly absorbed
into the surrogate's mean. For predicting damage, magnitude and timing are the
properties of the error that matter most.

\paragraph{Is $1\sigma$ a lot?} We estimate within-rater noise in three
published human studies \todo{cite the three studies in the third person} with a nearest-neighbour (semivariogram nugget)
estimator over repeated ratings by the same participant, and divide it by
$\sigma_f$, the standard deviation of that study's fitted objective over its
design box (Table~\ref{tab:anchor}). Real rating noise lands between $0.74$ and
$3.35$ landscape standard deviations, and participants are noisy from their
first rating, so on the onset-0 row of Table~\ref{tab:dose} these studies would
lose $12$--$20\%$ of the achievable improvement, interpolating log-linearly
between columns. The caveat is that $\sigma_f$ can only be computed through a fitted
oracle, the object this design exists to avoid, so this anchors the range of
interest rather than measuring it.

\begin{table}[t]
\centering
\caption{Three published human studies placed on the synthetic arm's axis.
$\sigma_f$ is the spread of the study's fitted objective over its design box;
the last column is the measured within-rater noise in those units.}
\label{tab:anchor}
\begin{tabular}{llrrr}
\toprule
study & ratings & $\sigma_f$ & noise / $\sigma_f$ & upper bound \\
\midrule
\texttt{ehmi} & as logged & 0.365 & 0.74 {\scriptsize $[0.48, 0.97]$} & 1.17 \\
\texttt{opticarvis} & as logged, two scales mixed & 0.374 & 3.35 {\scriptsize $[3.03, 3.58]$} & -- \\
\texttt{opticarvis} & one scale & 0.156 & 1.11 {\scriptsize $[0.81, 1.40]$} & 1.34 \\
\texttt{provoice} & as logged, one sign reversed & 0.298 & 1.02 {\scriptsize $[0.20, 1.56]$} (weak) & 2.58 \\
\texttt{provoice} & sign corrected & 0.449 & 0.60 {\scriptsize $[0.00, 0.90]$} (weak) & 2.37 \\
\bottomrule
\end{tabular}

\end{table}

\section{What makes an error big?}
\label{sec:currency}

Standardising every objective embeds an assumption: that a landscape's own
standard deviation is the right unit for its error. The rival is the improvement
actually available. A half-SD error is one thing when the optimum stands $0.76$
SDs above a random design and another when it stands $56.8$ above. We fit the
power law
\begin{equation}
  \mathbb{E}[R] \;=\; A\,\sigma_e^{\beta_c}\,\optz^{\beta_z}
\end{equation}
to the landscape-by-magnitude cell means of the raw excess regret $R$, using
quasi-Poisson estimating equations, which keep the cells in which noise happened
to help. If the spread is the right unit (\textsc{spread}), $\beta_z = 0$. If
what matters is the error relative to the achievable gain, then
$R = \optz\, h(\sigma_e/\optz)$, and for a power-law $h$ that means
$\beta_c + \beta_z = 1$ (\textsc{gain}). Intervals come from resampling
landscapes, the unit of replication for a landscape-level covariate.

\begin{table}[t]
\centering
\caption{Neither scale is right at either onset. Exponents of
$\mathbb{E}[R] = A\,\sigma_e^{\beta_c}\optz^{\beta_z}$ fitted to landscape
$\times$ magnitude cell means, with 95\% intervals from 2000 landscape
resamples. Each restriction is judged on the interval of the quantity it
constrains.}
\label{tab:currency}
% Eight columns overrun the text width at normal size. Set here rather than in
% tables/currency.tex, which is generated and would lose the fix on a rebuild.
\small
\setlength{\tabcolsep}{4pt}
\begin{tabular}{llrrcrcl}
\toprule
& & & \multicolumn{2}{c}{\textsc{spread}: $\beta_z = 0$} & \multicolumn{2}{c}{\textsc{gain}: $\beta_c + \beta_z = 1$} & \\
\cmidrule(lr){4-5} \cmidrule(lr){6-7}
error process & $t_0$ & $\beta_c$ & $\beta_z$ & 95\% CI & $\beta_c + \beta_z - 1$ & 95\% CI & verdict \\
\midrule
\textsc{gaussian} & 0 & 0.51 & 0.86 & [0.47, 1.13] & +0.37 & [+0.15, +0.70] & neither \\
\textsc{bias} & 0 & 0.51 & 0.85 & [0.49, 1.08] & +0.36 & [+0.17, +0.65] & neither \\
\textsc{drift} & 0 & 0.55 & 0.82 & [0.55, 1.11] & +0.37 & [$-$0.01, +0.63] & gain \\
\textsc{ar}(1) & 0 & 0.66 & 0.85 & [0.60, 1.15] & +0.51 & [+0.14, +0.77] & neither \\
\textsc{gaussian} & 20 & 0.62 & 1.15 & [0.64, 1.48] & +0.76 & [+0.03, +1.03] & neither \\
\textsc{bias} & 20 & 0.70 & 1.14 & [0.61, 1.58] & +0.85 & [+0.04, +1.08] & neither \\
\textsc{drift} & 20 & 0.65 & 1.15 & [0.61, 1.55] & +0.80 & [$-$0.00, +1.14] & gain \\
\textsc{ar}(1) & 20 & 0.71 & 1.14 & [0.64, 1.49] & +0.85 & [+0.13, +1.11] & neither \\
\bottomrule
\end{tabular}

\end{table}

Neither restriction holds at either onset (Table~\ref{tab:currency}). $\beta_z$
excludes zero in all eight conditions, so standardising an objective does not
make error magnitudes comparable across problems, and a noise level reported
without the problem's headroom is under-specified. The cost also grows with
$\optz$ faster than \textsc{gain} allows: $\beta_c + \beta_z - 1$ is $+0.36$ to
$+0.51$ from the first observation and $+0.76$ to $+0.85$ from iteration 21,
every interval excluding zero, and OLS on the logarithm of the positive cells
returns the same verdicts. Read descriptively, $\beta_z$ is $0.82$ to $0.86$
from the first observation, with intervals that include one, so a given $\sigma_e$ destroys a roughly constant \emph{fraction} of the
achievable improvement; from iteration 21 it is $1.14$ to $1.15$. Appendix~\ref{sec:manipulation}
shows that $\optz$ is not what does the work: raising a narrow optimum changes
nothing, and the dependence runs through how findable the optimum is.

\section{The cost is mediated by a one-shot selection loss}
\label{sec:mediation}

Both preceding sections describe the cost of error in terms of properties of the
landscape, but neither gives a practitioner a number. Define
$\frag(\sigma_e)$ as the expected loss from a single greedy choice among $m$
quasi-random candidates when each is observed with $\mathcal{N}(0,\sigma_e^2)$
error:
\begin{equation}
  \frag(\sigma_e) \;=\;
  \mathbb{E}_{\epsilon}\Big[\max_i g_i - g_{\arg\max_i (g_i + \epsilon_i)}\Big],
  \qquad \epsilon_i \sim \mathcal{N}(0,\sigma_e^2),
  \label{eq:frag}
\end{equation}
where $g_i = g(x_i)$ over a fixed Sobol design. It involves no surrogate, no
sequence and no acquisition function; at $m = 256$ and $4000$ Monte-Carlo draws
it costs about $0.2$ seconds per landscape. Being indexed by the error level as
well as the landscape, it varies along both axes of the sweep.

\begin{table}[t]
\centering
\caption{Mediation. Fitted on landscape $\times$ condition cell means,
predictors $z$-scored, every model also carrying onset and error-process
indicators; intervals and $p$-values from a 2000-replicate bootstrap over
landscapes. Conditional on the one-shot selection loss, the error magnitude
explains nothing.}
\label{tab:mediation}
% Seven numeric columns with long math headers run 22pt past the text width at
% normal size. Set here rather than in tables/mediation.tex, which is generated
% and would lose the fix on the next rebuild.
\small
\setlength{\tabcolsep}{5pt}
\begin{tabular}{lrrrrrrrr}
\toprule
model & $\mathrm{frag}(\sigma_e)$ & $\log\sigma_e$ & $\log\mathrm{opt}_z$ & $\log\sigma_e{\times}\log\mathrm{opt}_z$ & $\log$ tail & rugged. & $d$ & $R^2$ \\
\midrule
indicators only & -- & -- & -- & -- & -- & -- & -- & 0.040 \\
$\log\sigma_e$ only & -- & +0.37$^{***}$ & -- & -- & -- & -- & -- & 0.158 \\
$\mathrm{frag}$ only & +0.75$^{***}$ & -- & -- & -- & -- & -- & -- & 0.520 \\
descriptors only & -- & +0.37$^{***}$ & +0.64$^{***}$ & -- & +0.20$^{**}$ & $-$0.09 & $-$0.03 & 0.526 \\
both & +0.66$^{***}$ & $-$0.07 & +0.42$^{**}$ & -- & +0.20$^{**}$ & $-$0.06 & $-$0.03 & 0.701 \\
descriptors + int. & -- & +0.37$^{***}$ & +0.64$^{***}$ & +0.54$^{***}$ & +0.20$^{**}$ & $-$0.09 & $-$0.03 & 0.779 \\
both + int. & +0.31$^{***}$ & +0.17$^{*}$ & +0.54$^{**}$ & +0.40$^{***}$ & +0.20$^{**}$ & $-$0.08 & $-$0.03 & 0.792 \\
\bottomrule
\end{tabular}

\end{table}

With onset and error-process indicators in every model, $\frag(\sigma_e)$
explains $R^2 = 0.52$ of the variation in excess regret across
landscape-by-condition cells ($0.48$ on its own), nearly as much as the primary
landscape descriptors together with the error magnitude ($R^2 = 0.53$;
Table~\ref{tab:mediation}). Entered
jointly the two are complementary ($R^2 = 0.70$), and the informative part is
which coefficient survives: $\frag$ keeps $\beta = +0.66$ with a bootstrap
interval of $[0.28, 0.74]$, while the error magnitude collapses to
$\beta = -0.07$ with $p = 0.44$. \textbf{Given what one noisy choice would cost on this landscape, the error
magnitude adds no detectable information about what fifty will cost.} Dimension and
ruggedness are null, and holding $\optz$ fixed confirms that dimension does
little to excess regret (Appendix~\ref{sec:manipulation}).

Equation~\ref{eq:frag} needs only a design space and an estimate of rater noise,
both of which a practitioner has before running a study. The residual terms say where the one-shot approximation breaks down: $\optz$
($+0.42$) and tail weight ($+0.20$) keep significant positive coefficients, so
on tall, heavy-tailed landscapes fifty sequential noisy choices cost more than
one would predict from a single one.

\section{Which acquisition function?}
\label{sec:acq}

\begin{table}[t]
\centering
\caption{Acquisition robustness, ranked over landscapes within each condition
and averaged over conditions. ``Absolute loss'' is the deployment number: the
post-onset regret under noise as a percentage of the achievable gain.}
\label{tab:acq}
\begin{tabular}{lrrr}
\toprule
acquisition & mean rank & excess regret & absolute loss \\
\midrule
qUCB & 3.95 & 4.14\% & 27.9\% \\
qNEI & 4.34 & 4.89\% & 25.2\% \\
UCB & 4.41 & 4.73\% & 27.4\% \\
LogEI & 4.61 & 5.21\% & 26.0\% \\
qEI & 4.73 & 5.37\% & 25.5\% \\
EI & 4.98 & 5.42\% & 26.2\% \\
qPI & 6.29 & 7.90\% & 31.0\% \\
Greedy & 6.56 & 8.11\% & 32.5\% \\
LogPI & 7.48 & 9.19\% & 31.3\% \\
PI & 7.66 & 9.44\% & 31.7\% \\
\midrule
model-free floor & -- & 0.00\% & 41.0\% \\
\bottomrule
\end{tabular}

\end{table}

The confidence-bound family and noisy expected improvement are the most robust,
and the probability-of-improvement family and pure exploitation (Greedy) fill
the bottom four places on both robustness and absolute loss
(Table~\ref{tab:acq}). Every model-based arm is far better than the model-free
floor at $41.0\%$, so no arm wins by declining to learn. The Friedman test over the twenty landscapes as blocks reaches
$p < 0.05$ in $28$ of $32$ conditions, and $115$ of $288$ pairwise comparisons
against the per-condition winner survive Benjamini--Hochberg correction
\citep{benjamini1995controlling}.

The landscapes also disagree. The median Kendall's $W$ over the 32 conditions
is $0.247$: the landscapes largely do not agree on an ordering. qUCB
takes $10$ of $32$ conditions, UCB and qNEI $6$ each, LogEI $5$, EI $3$, qEI
$2$. A study that measured any one of these landscapes and reported its winner
would have produced a defensible recommendation that largely does not
transfer, and a study run on one design space cannot tell whether its own does. Appendix~\ref{sec:incumbent} shows that part of this ranking is
not about the acquisitions at all.

\section{What the fitted-oracle design would have concluded}
\label{sec:fitted}

The introduction argued that fitting a regression oracle to archival ratings
entangles the cost of feedback error with the cost of the oracle being wrong. To
measure this rather than argue it, we run the data-driven simulator on the three
archival datasets at the same error magnitudes, expressed in each dataset's own
$\sigma_f$ so that ``$1\sigma$'' means the same thing in both arms.

\begin{table}[h]
\centering
\caption{The same experiment, run against exact functions and against oracles
fitted to ratings, under \textsc{gaussian} error. Percentages are post-onset
excess as a fraction of the gap between the optimum and a same-budget model-free
floor (the floor gap), the one normaliser computable in both arms. The fitted
design reports roughly a third of the cost.}
\label{tab:fitted}
\begin{tabular}{lrrrrrrrr}
\toprule
& \multicolumn{4}{c}{error from it.\ 1} & \multicolumn{4}{c}{error from it.\ 21} \\
\cmidrule(lr){2-5} \cmidrule(lr){6-9}
arm & $0.05\sigma$ & $0.25\sigma$ & $1\sigma$ & $5\sigma$ & $0.05\sigma$ & $0.25\sigma$ & $1\sigma$ & $5\sigma$ \\
\midrule
known functions (20) & 10 & 44 & 99 & 148 & 3 & 7 & 10 & 16 \\
fitted oracle (3) & 1 & 11 & 33 & 47 & 1 & 3 & 6 & 8 \\
fitted, no augmentation (3) & 2 & 14 & 38 & 50 & 1 & 4 & 6 & 7 \\
\bottomrule
\end{tabular}

\end{table}

The fitted arm has no verified optimum, so Table~\ref{tab:fitted}'s denominator
uses the best value any arm reached. That value is attainable by construction,
which shrinks the denominator and inflates the fitted arm's reported fraction,
so this bias works against the result below. With three design spaces against
twenty, the fitted intervals are also wide.

\textbf{The fitted design understates the cost of feedback error by roughly
threefold.} At $1\sigma$ from the first rating, the known-function arm puts the
loss at $99\%$ of the floor gap (cluster-bootstrap interval $[55, 155]$) and the
fitted arm at $33\%$ ($[5, 75]$). At $0.05\sigma$ the fitted arm reports $1\%$
($[-4, 3]$), indistinguishable from no effect, where the exact functions report
$10\%$ ($[4, 18]$). The direction is the same at all four magnitudes and both
onsets; the intervals overlap at the larger magnitudes, so this is a consistent
gap rather than a separation in every cell.

\textbf{Oracle fidelity does not explain it.} Turning off the deployed jitter
augmentation improves the companion's oracles by up to $0.29$ held-out $R^2$
(\texttt{opticarvis}, $0.52$ to $0.81$). Rerunning the companion arm without it
changes the fraction destroyed by $+0.016$ (median $+0.006$) over the $2{,}400$
cells the two arms share, and the augmentation-free arm is the worse of the two
in $52.6\%$ of them. If the gap were fidelity, it would close as the oracle
improved. What the fitted design changes is what the regret is measured
\emph{against}: a random-search estimate of the optimum that the optimizer
reaches routinely, which compresses the range in which any error can be seen to
do damage. The onset effect compresses with it, to $0.8$--$5.6\times$ across
magnitudes against $3.6$--$9.6\times$ in the exact arm.

\section{When the error is in the design, not the rating}
\label{sec:inputerror}

Every error process above corrupts the reported value. People also make the
other mistake: a slider overshoots, the wrong option is pressed, a configuration
is applied to the wrong condition. The optimizer proposes $x$, the person
experiences $x' \neq x$, rates it honestly, and the rating is filed against $x$:
the right value at the wrong location. This is also the corruption a
fitted-oracle design is least able to measure, because scoring it needs the
objective at $x'$, where a fitted oracle adds an error of its own.

A \textsc{slip} displaces every trial by $\delta \sim \mathcal{N}(0, s^2)$ per
coordinate, with $s$ a fraction of that coordinate's range, and clips the result
to the box, as a bounded control stops at its end. A \textsc{misclick} leaves the
trial exact with probability $1 - p$ and otherwise lands on a uniformly random
design. Both magnitudes are swept over $\{0.01, 0.05, 0.15, 0.4\}$ at the same
onsets and budget as before, with six model-based acquisitions (LogEI, EI, PI,
UCB, qUCB, qNEI), the two floors and five seeds. Neither magnitude is in
landscape standard deviations, so this section's numbers can be compared with
the response-error grid only on the damage axis. The rating itself is exact.

If the slip goes unnoticed, the surrogate trains on $(x, g(x'))$ and the design
eventually deployed is the recorded $x$, so we score the deployed design there. A counterfactual arm records
$(x', g(x'))$ instead, as an instrumented interface would, with identical slip
draws; the contrast separates the cost of evaluating the wrong point from the
cost of mislabelling it. The model-free floors are no longer a zero control:
they never read an observation, but they do evaluate the wrong points, so their
excess regret is the geometric cost of a slip with no learning to corrupt.

\begin{table}[t]
\centering
\caption{Fraction of the achievable improvement destroyed by an input error. A
slip's magnitude is the displacement SD as a fraction of each coordinate's
range; a misclick's is the probability of a random trial. The last three
columns decompose the early-onset slip: \emph{floor}, the model-free arms;
\emph{logged}, a learner whose record holds the point evaluated;
\emph{mislabelling}, the paired difference between the unnoticed slip (second
column) and \emph{logged}. Brackets: 95\% intervals over landscapes.}
\label{tab:inputerror}
\small\setlength{\tabcolsep}{3pt}
\begin{tabular}{lrrrrrrr}
\toprule
& \multicolumn{2}{c}{\textsc{slip}} & \multicolumn{2}{c}{\textsc{misclick}} & \multicolumn{3}{c}{\textsc{slip} from it.\ 1, decomposed} \\
\cmidrule(lr){2-3} \cmidrule(lr){4-5} \cmidrule(lr){6-8}
magnitude & \multicolumn{1}{c}{from it.\ 1} & \multicolumn{1}{c}{from it.\ 21} & \multicolumn{1}{c}{from it.\ 1} & \multicolumn{1}{c}{from it.\ 21} & \multicolumn{1}{c}{floor} & \multicolumn{1}{c}{logged} & \multicolumn{1}{c}{mislabelling} \\
\midrule
1\% & 1.5 \makebox[3.0em][l]{\scriptsize [0, 3]} & 1.1 \makebox[3.0em][l]{\scriptsize [0, 2]} & 0.4 \makebox[3.0em][l]{\scriptsize [0, 1]} & 0.1 \makebox[3.0em][l]{\scriptsize [0, 0]} & 0.3 \makebox[3.0em][l]{\scriptsize [0, 1]} & 0.7 \makebox[3.0em][l]{\scriptsize [0, 2]} & 0.8 \makebox[3.0em][l]{\scriptsize [0, 2]} \\
5\% & 5.4 \makebox[3.0em][l]{\scriptsize [3, 8]} & 3.0 \makebox[3.0em][l]{\scriptsize [2, 5]} & 2.8 \makebox[3.0em][l]{\scriptsize [2, 4]} & 0.5 \makebox[3.0em][l]{\scriptsize [0, 1]} & 0.2 \makebox[3.0em][l]{\scriptsize [$-$1, 1]} & 2.8 \makebox[3.0em][l]{\scriptsize [2, 4]} & 2.7 \makebox[3.0em][l]{\scriptsize [1, 4]} \\
15\% & 11.5 \makebox[3.0em][l]{\scriptsize [8, 15]} & 5.0 \makebox[3.0em][l]{\scriptsize [3, 8]} & 6.8 \makebox[3.0em][l]{\scriptsize [5, 9]} & 2.0 \makebox[3.0em][l]{\scriptsize [1, 3]} & 1.2 \makebox[3.0em][l]{\scriptsize [$-$1, 3]} & 7.1 \makebox[3.0em][l]{\scriptsize [5, 10]} & 4.4 \makebox[3.0em][l]{\scriptsize [2, 6]} \\
40\% & 23.2 \makebox[3.0em][l]{\scriptsize [16, 31]} & 6.4 \makebox[3.0em][l]{\scriptsize [4, 10]} & 13.5 \makebox[3.0em][l]{\scriptsize [10, 18]} & 3.5 \makebox[3.0em][l]{\scriptsize [2, 6]} & 6.4 \makebox[3.0em][l]{\scriptsize [0, 13]} & 19.6 \makebox[3.0em][l]{\scriptsize [12, 28]} & 3.6 \makebox[3.0em][l]{\scriptsize [2, 6]} \\
\bottomrule
\end{tabular}

\end{table}

\textbf{A slip costs about as much as a rating error.} A slip of $5\%$ of each
coordinate's range on every trial, from the first, destroys $5.4\%$ of the
achievable improvement; one of $15\%$ destroys $11.5\%$; one of $40\%$, $23.2\%$
(Table~\ref{tab:inputerror}). On the damage axis these sit beside the
$0.25\sigma$, $1\sigma$ and $5\sigma$ rating errors of Table~\ref{tab:dose}, at
$5.3\%$, $14.1\%$ and $22.0\%$. A misclick is cheaper: even when $40\%$ of
trials land somewhere random it costs $13.5\%$, less than a $15\%$ slip on every
trial. A late slip is again cheaper than an early one, but by less than a late
rating error: the early-to-late ratio is $1.4$--$3.6\times$ against $5.0\times$
at $1\sigma$ in Table~\ref{tab:dose}.

\textbf{Half of a small slip's cost is the record, not the displacement.}
The decomposition in Table~\ref{tab:inputerror} (in full,
Table~\ref{tab:inputmech}) separates the mechanisms. The model-free floor barely
moves, since a random design displaced is still a random design; only at $40\%$,
where clipping piles evaluations at the edges of the box, does it lose $6.4\%$
$[0, 13]$. A learner whose record is correct loses far more, $2.8\%$ at a $5\%$
slip and $19.6\%$ at $40\%$: a proposal displaced is a worse proposal.
Mislabelling adds $2.7$ points at $5\%$ and $4.4$ at $15\%$, half and two-fifths
of the total; at $40\%$ the displacement dominates and the mislabelled share
falls to a sixth. An interface that records what was applied recovers that
share at no cost in trials.

\textbf{The design in the log is not the design that earned its rating.} Under
an unnoticed slip the practitioner deploys the recorded $x$. Scored there, the
final design's excess regret is $16.4\%$ of the achievable improvement at a
$5\%$ slip from the first trial, $34.3\%$ at $15\%$ and $65.8\%$ at $40\%$,
against $6.5\%$, $13.4\%$ and $24.9\%$ for the best design the person actually
experienced (Appendix~\ref{app:inputerror}): a factor of $1.6$ to $2.6$. The
optimizer's progress survives an input error far better than the record of it
does, and the record is what a deployment inherits.

\section{The conclusions under six further checks}
\label{sec:robustness}

Each check changes one thing about the design and is reported in full in an
appendix.

\paragraph{The surrogate's noise estimate is not the problem.} At $1\sigma$ and
above the GP's fitted noise is an eighth of the injected value for the first
twenty iterations, yet
telling it the true variance removes only $5.0\%$ of the cost, and no condition
survives correction. The cost is an information cost, and the onset effect is not
a failure of noise estimation (Appendix~\ref{sec:noisediag}).

\paragraph{The incumbent is part of the method.} Switching improvement-based
acquisitions from the posterior-mean incumbent to the observed maximum removes
$15.6\%$ of the cost overall and $35$--$41\%$ for the probability-of-improvement
family, so a robustness ranking over acquisitions ranks acquisition-and-incumbent
pairs (Appendix~\ref{sec:incumbent}).

\paragraph{Most of the apparent dimension effect is signal strength.} Landscapes
built to vary one property at a time show that $\optz$ is not causal on its own,
and holding $\optz$ fixed removes most of the dimension effect that a Levy ladder
appeared to show (Appendix~\ref{sec:manipulation}).

\paragraph{Multi-objective BO is no more fragile.} On six BoTorch
problems with published hypervolume optima the dose-response and onset
structures replicate, early error costs roughly half as much as on the scalar
suite on a common normaliser, and the four hypervolume acquisitions are
indistinguishable (Appendix~\ref{sec:multiobjective}).

\paragraph{The budget matters; the rating scale cuts both ways.} With the late
onset held at $40\%$ of the run, the early-to-late cost ratio is $2.3\times$,
$5.0\times$ and $11.7\times$ at $T = 25$, $50$ and $100$, so every onset number
above describes a fifty-iteration study. A bounded, discretised rating scale
makes small errors costlier and large ones slightly cheaper, and qKG and
replication, two methods built for noise, recover about a third of the
early-error cost (Appendix~\ref{sec:designchoices}).

\paragraph{The replication is void; post hoc, the claims hold.} We
preregistered four claims before running ten fresh seeds. One control failed,
a landscape falling below the preregistered headroom threshold, so the run is
void by its own rule; with that landscape excluded post hoc, all four claims
clear their thresholds (Appendix~\ref{sec:confirmatory}).

\section{Discussion}
\label{sec:discussion}

At measured rater-noise levels, feedback error costs a study $12$--$20\%$ of the
improvement it could have had: worth reducing, but not a catastrophe.
Classifying \emph{what kind} of error one's raters make is largely wasted
effort, since at $1\sigma$ and above the four processes are within $32\%$ of
each other. Choosing an acquisition function is worth about $2.8\times$ in cost
at $1\sigma$, but the ranking transfers only as far as a new design space
resembles the one it was measured on, which at a median $W$ of $0.247$ is not
far. What does transfer is Equation~\ref{eq:frag}: computed for the design space
and the rater noise, it predicts half the variation in damage ($R^2 = 0.48$ on
its own). When the error is in the design, the interface matters as much as the
optimizer: recording what was applied removes up to half of a small slip's
cost, and the design a log recommends should be re-evaluated before it ships.

Methodologically, the fitted-oracle design recovers the same qualitative
structure at about a third of the magnitude, and a better oracle does not close
the gap (Section~\ref{sec:fitted}). A study that measures robustness against a
random-search estimate of its own optimum measures a quantity whose ceiling it
set itself, a mistake that is cheap to avoid.

Two things remain unexplained. The onset effect is large, is not a failure of
noise estimation (Appendix~\ref{sec:noisediag}), and grows with the budget
(Appendix~\ref{sec:designchoices}); we have no mechanism for it. And the cost of
error grows with $\optz$ faster than \textsc{gain} allows
(Section~\ref{sec:currency}), yet $\optz$ is not causal on its own
(Appendix~\ref{sec:manipulation}): an interaction between signal strength and
findability, characterised but not modelled.

\paragraph{Limitations.} The landscapes are analytic, not people. They remove
the surrogate-human confound at the cost of the realism the fitted-oracle design
buys, and Section~\ref{sec:fitted} shows that the design changes the answer but
cannot say which answer is closer to a person. The rater-noise anchor runs through a fitted oracle's $\sigma_f$, and every
onset number describes a fifty-iteration study. The multi-objective arm is six problems at
five seeds; it cannot separate the hypervolume acquisitions, and the mediator,
which has no Pareto analogue, is untested there. The input-error arm covers six
acquisitions at five seeds, one slip distribution and one misclick model. The
robust baselines are two, not a survey. And the preregistered replication is
void on a control failure, so the claims clear their thresholds only after a
post-hoc exclusion and cannot be called confirmed.

\section*{Reproducibility statement}
The full pipeline is scripted end to end, from the landscape statistics through
the sweep to every table in this paper, and every stochastic component is
explicitly seeded. Regenerating published runs from their own recorded settings
reproduces every model-based run bit for bit. The one exception is the
model-free floors of the main sweep, whose noise seed moved when two robust
baselines were added to the acquisition list after that sweep ran: their
candidates and regret still reproduce exactly, since a floor never reads an
observation, but their recorded noisy observations do not. No reported number
depends on those, and the seed order is now pinned by test. Three further checks are automated. The vendored
benchmark implementations are asserted by test to agree with their source
implementations to $10^{-9}$ relative on random points, so the objective cannot
silently drift. The model-free control of Section~\ref{sec:controls} is verified
to be exactly zero in every run. And one recorded optimum in the source
suite, the supremum for a superseded parameter setting, is wrong by $64\%$; we
detected this by searching every box against its recorded value and corrected it
here; the correction and the provenance of the original are
both asserted by test. \todo{repository URL and commit at camera-ready.}

% Required by the ICLR 2026 LLM policy (iclr.cc/FAQ/LLM): significant LLM use
% must be described in its own section, which does not count toward the limit.
\section*{LLM usage}
A large language model (Claude, from Anthropic) was used extensively in this
work. It proposed and implemented parts of the experimental design and analysis,
wrote and tested most of the simulation and analysis code, ran and monitored the
experiments, drafted and edited the manuscript, and checked the numbers in the
text against the analysis outputs. The authors set the research questions, chose
among the proposed designs, and are responsible for everything reported here.
\todo{authors: confirm or amend this description before submission.}

\bibliographystyle{iclr2026_conference}
\bibliography{references}

\clearpage
\appendix
\section{Ablations}
\suppressfloats[t]

\subsection{Does the surrogate know it is being lied to?}
\label{sec:noisediag}
The GP fits its observation noise as a free hyperparameter under a prior that
shrinks it, so a measured cost of noise could in principle be the surrogate's
failure to notice the error rather than the information the error destroys. This
is the first thing to rule out, because if it were the explanation then every
number above would be a statement about BoTorch's default prior rather than
about noisy feedback.

The mis-estimation is real, and it lands where it would do the most damage
(Table~\ref{tab:noisediag}). Refitting each surrogate from the sweep's own logs
using only its first $n$ observations, at $1\sigma$ and above the fitted noise standard
deviation is about an eighth of the injected one for the first twenty of fifty
iterations and has converged by the end of the run, while at $0.05\sigma$ it is
overestimated by a quarter to a half throughout --- so a diagnostic evaluated only at
$n = 50$ looks reassuring and is misleading. Refitting with a wide noise prior
recovers the injected level far sooner ($0.55$ rather than $0.14$ at $n = 15$,
$1\sigma$), so the shrinkage is the prior's doing and not a limit of the data.

\begin{table}[t]
\centering
\caption{Fitted over injected observation-noise standard deviation, as a
function of how many observations the surrogate has seen, for error present from
iteration~1. The estimate is badly wrong exactly while the trajectory is being
decided --- and Table~\ref{tab:knownnoise} shows that this costs almost nothing.}
\label{tab:noisediag}
\begin{tabular}{lrrrr}
\toprule
observations & $0.05\sigma$ & $0.25\sigma$ & $1\sigma$ & $5\sigma$ \\
\midrule
8 & 1.43 & 0.30 & 0.12 & 0.11 \\
15 & 1.50 & 0.41 & 0.14 & 0.10 \\
25 & 1.37 & 0.72 & 0.66 & 0.56 \\
50 & 1.24 & 0.88 & 0.83 & 0.89 \\
\bottomrule
\end{tabular}

\end{table}

It was natural to expect this to explain the onset effect of
Section~\ref{sec:cost}: from the first observation the error arrives in exactly
the window where the noise cannot yet be estimated, whereas at a mid-run onset
the surrogate has a settled length scale and thirty iterations in which to
learn. \textbf{We tested that directly and it is wrong.} Supplying the true
injected variance as \texttt{train\_Yvar} and re-running $4{,}800$ paired cells
removes only $5.0\%$ of the measured cost, no condition survives FDR correction
(Table~\ref{tab:knownnoise}), and the onset-0 to onset-20 ratio at $1\sigma$ does
not shrink --- it goes from $4.0$ to $5.2$. Per landscape the largest reduction
is Ackley at $47\%$; per acquisition the effect is inconsistent in sign, helping
PI ($22\%$) and qNEI ($15\%$) but \emph{hurting} qUCB ($-29\%$).

Two conclusions follow. Bayesian optimization is remarkably insensitive to its
own noise hyperparameter in this regime, despite getting it badly wrong for the
first two-fifths of the budget --- worth knowing on its own, and a caution
against reading too much into noise-estimation diagnostics. And the cost this
paper measures is the information cost of the error, not an artefact of a
mis-specified surrogate, which is what licenses reading Table~\ref{tab:dose} as
a property of the problem. The onset effect itself is therefore left
unexplained; we return to it in Section~\ref{sec:discussion}.

\begin{table}[t]
\centering
\caption{Telling the surrogate the truth about its noise. Paired within
(landscape, acquisition, magnitude, onset, seed); positive $\Delta$ means the
treatment suffers more. Nothing survives FDR correction.}
\label{tab:knownnoise}
\begin{tabular}{rrrrrrrr}
\toprule
$\sigma_e$ & $t_0$ & reference & treatment & $\Delta$ & share & $d_z$ & $p_{\text{FDR}}$ \\
\midrule
0.05 & 0 & 0.065 & 0.077 & +0.012 & $-$19\% & 0.03 & 0.888 \\
0.25 & 0 & 0.271 & 0.166 & $-$0.105 & +39\% & $-$0.18 & 0.888 \\
1 & 0 & 0.790 & 0.848 & +0.058 & $-$7\% & 0.10 & 0.888 \\
5 & 0 & 1.484 & 1.463 & $-$0.021 & +1\% & $-$0.04 & 0.888 \\
0.05 & 20 & 0.002 & 0.001 & $-$0.002 & +71\% & $-$0.02 & 0.913 \\
0.25 & 20 & 0.072 & 0.041 & $-$0.031 & +43\% & $-$0.20 & 0.888 \\
1 & 20 & 0.197 & 0.164 & $-$0.034 & +17\% & $-$0.19 & 0.888 \\
5 & 20 & 0.531 & 0.483 & $-$0.048 & +9\% & $-$0.16 & 0.888 \\
\bottomrule
\end{tabular}

\end{table}

\FloatBarrier
\subsection{Is it the acquisition or the incumbent?}
\label{sec:incumbent}
$\mathrm{best}_f$ is the maximum posterior mean at visited points, which shrinks
under noise and so makes improvement-based acquisitions automatically more
exploratory. The acquisitions differ in their exposure to this: LogEI, EI, PI,
LogPI and qEI consume $\mathrm{best}_f$; UCB does not; qNEI ignores it entirely.
A ranking of acquisitions by robustness is therefore partly a ranking of their
sensitivity to a design choice that is not part of the acquisition at all.
Re-running with the observed maximum as the incumbent removes $15.6\%$ of the
measured cost overall, with four of eight conditions significant after FDR
correction (Table~\ref{tab:incumbent}).

The per-acquisition breakdown is in Table~\ref{tab:incumbentacq}. UCB
and qNEI move by \emph{exactly} zero, as they must --- a control that validates
the pairing at machine precision. The probability-of-improvement family moves
most: LogPI loses $41\%$ of its measured cost and PI $35\%$, against $17\%$ for
EI, $2\%$ for LogEI and $-2\%$ for qEI. So a substantial part of ``PI is fragile
under noise'' is really ``the posterior-mean incumbent is fragile under noise,
and PI is the acquisition most exposed to it''. Every acquisition-robustness
ranking here, including Table~\ref{tab:acq}, is a ranking of
acquisition-and-incumbent pairs, and the same is true of the applied results it
is modelled on. Reporting the incumbent definition alongside the acquisition
should be standard.

\begin{table}[t]
\centering
\caption{Observed-maximum against posterior-mean incumbent, same pairing as
Table~\ref{tab:knownnoise}.}
\label{tab:incumbent}
\begin{tabular}{rrrrrrrr}
\toprule
$\sigma_e$ & $t_0$ & reference & treatment & $\Delta$ & share & $d_z$ & $p_{\text{FDR}}$ \\
\midrule
0.05 & 0 & 0.010 & 0.003 & $-$0.007 & +72\% & $-$0.56 & 0.019 \\
0.25 & 0 & 0.050 & 0.026 & $-$0.024 & +48\% & $-$1.39 & $<$0.001 \\
1 & 0 & 0.132 & 0.100 & $-$0.032 & +24\% & $-$0.93 & 0.004 \\
5 & 0 & 0.223 & 0.200 & $-$0.022 & +10\% & $-$0.71 & 0.011 \\
0.05 & 20 & 0.003 & 0.001 & $-$0.002 & +81\% & $-$0.24 & 0.142 \\
0.25 & 20 & 0.014 & 0.005 & $-$0.009 & +64\% & $-$0.88 & 0.002 \\
1 & 20 & 0.027 & 0.021 & $-$0.005 & +20\% & $-$0.50 & 0.032 \\
5 & 20 & 0.049 & 0.048 & $-$0.000 & +0\% & $-$0.01 & 0.421 \\
\bottomrule
\end{tabular}

\end{table}

\begin{table}[t]
\centering
\caption{The incumbent arm by acquisition. UCB and qNEI never read
$\mathrm{best}_f$ and move by exactly zero, which validates the pairing; the
probability-of-improvement family moves most.}
\label{tab:incumbentacq}
\begin{tabular}{lrrrr}
\toprule
acquisition & reference & treatment & $\Delta$ & share removed \\
\midrule
LogPI & 0.555 & 0.328 & $-$0.227 & +41\% \\
PI & 0.572 & 0.373 & $-$0.199 & +35\% \\
EI & 0.459 & 0.382 & $-$0.077 & +17\% \\
LogEI & 0.493 & 0.481 & $-$0.011 & +2\% \\
qNEI & 0.378 & 0.378 & +0.000 & +0\% \\
UCB & 0.346 & 0.346 & +0.000 & +0\% \\
qEI & 0.445 & 0.454 & +0.009 & $-$2\% \\
\bottomrule
\end{tabular}

\end{table}

\FloatBarrier
\subsection{Breaking the descriptor confound}
\label{sec:manipulation}
In our suite $\optz$, sparsity and skew correlate at about $0.9$, and for a
bounded function with an isolated optimum that is close to a structural
identity, so Section~\ref{sec:mediation}'s descriptor terms cannot be
separated observationally, however many real benchmarks are added. Two
purpose-built families vary them by construction: a cos-field background plus one
spike, whose amplitude moves $\optz$ and whose width moves sparsity; and the Levy
function at $d \in \{4, 7, 11\}$, one landscape family across a dimension ladder
(Table~\ref{tab:manipulation}).

\textbf{$\optz$ on its own is not causal.} Holding the spike narrow and raising
its amplitude by $4\times$ moves $\optz$ from $5.0$ to $21.7$ and changes the
excess regret at $1\sigma$ from $0.246$ to $0.245$ --- no effect at all. Widen
the spike and the same amplitude change raises excess from $0.495$ to $1.888$.
What the observational $\beta_{\optz} = +0.42$ picks up is therefore an
interaction, not a main effect: a taller optimum costs more to miss only when it
is broad enough that the optimizer would otherwise have found it. On a true
needle, the optimizer was never going to find the optimum with or without noise,
so there is nothing for the noise to destroy. It also means that the $\optz$ exponent of Section~\ref{sec:currency}
describes the suite rather than a mechanism: it is not $\optz$ as such that
makes an error expensive.

\textbf{Dimension, and a within-family ladder that corrects itself.} Within the
Levy family, excess regret at $1\sigma$ rises monotonically from $0.114$ at
$d = 4$ to $0.222$ at $d = 7$ and $0.264$ at $d = 11$, while the cross-suite
regression puts $\beta_d$ at $-0.03$ and not significant. Read on its own that
looks like a real effect the observational analysis missed. It is not, or not
mostly: $\optz$ also rises along the Levy ladder, $1.53$ to $2.07$ to $2.62$, so
that comparison varies dimension and signal strength together.

The \texttt{bump} ladder separates them. Amplitude and width are chosen at each
dimension so that $\optz$ is $9.0$ and the spike occupies a $10^{-6}$ fraction of
the box at all three, leaving dimension as the only difference. Excess regret is
then $0.031$ at $d = 4$, $0.028$ at $d = 7$ and $0.048$ at $d = 11$ --- not
monotone, and a factor of $1.5$ across the range where the confounded ladder gave
$2.3$. \textbf{Most of what looked like a dimension effect was $\optz$.} What
does rise monotonically with dimension is the inference metric, the true value of
the design a practitioner would actually deploy: $0.52$, $0.70$, $0.99$. Higher
dimension does not destroy more of the optimizer's progress; it makes the
progress harder to identify from noisy observations at the end.

The methodological lesson is that a landscape-level regression over a
hand-assembled suite can miss an effect and a within-family ladder is the cheap
fix, but a family indexed by dimension is not automatically a manipulation of
dimension; ours became one only when it was built to hold $\optz$ fixed.

\begin{table}[t]
\centering
\caption{Purpose-built landscapes. Raw post-onset excess regret in landscape
standard deviations, \textsc{gaussian} error from iteration~1. Amplitude and width are
varied independently in the top block; the bottom block is one function family
at three dimensions.}
\label{tab:manipulation}
\begin{tabular}{lrrrrrrr}
\toprule
landscape & $d$ & $\mathrm{opt}_z$ & sparsity & $0.05\sigma$ & $0.25\sigma$ & $1\sigma$ & $5\sigma$ \\
\midrule
\texttt{bump\_a4\_w0.05} & 4 & 5.0 & $<$1.5e-05 & 0.007 & 0.152 & 0.246 & 0.324 \\
\texttt{bump\_a16\_w0.05} & 4 & 21.7 & $<$1.5e-05 & 0.014 & 0.026 & 0.245 & 0.262 \\
\addlinespace
\texttt{bump\_a4\_w0.15} & 4 & 5.3 & 3.8e-04 & 0.133 & 0.289 & 0.495 & 0.704 \\
\texttt{bump\_a16\_w0.15} & 4 & 12.3 & 1.8e-04 & 0.298 & 0.730 & 1.888 & 3.173 \\
\midrule
\texttt{levy\_4d} & 4 & 1.5 & 2.6e-01 & 0.006 & 0.048 & 0.114 & 0.200 \\
\texttt{levy\_7d} & 7 & 2.1 & 7.3e-02 & 0.051 & 0.115 & 0.222 & 0.305 \\
\texttt{levy\_10} & 11 & 2.6 & 1.3e-02 & 0.001 & 0.147 & 0.264 & 0.481 \\
\midrule
\texttt{bump\_d4} & 4 & 9.0 & $<$1.5e-05 & 0.001 & 0.121 & 0.279 & 0.327 \\
\texttt{bump\_d7} & 7 & 9.0 & $<$1.5e-05 & $-$0.001 & 0.103 & 0.248 & 0.428 \\
\texttt{bump\_d11} & 11 & 9.0 & $<$1.5e-05 & 0.142 & 0.215 & 0.427 & 0.734 \\
\bottomrule
\end{tabular}

\end{table}

\FloatBarrier
\section{Does any of this hold for multiple objectives?}
\label{sec:multiobjective}

The applied setting is frequently multi-objective \citep{daulton2021parallel}
and everything above is scalar, so the generalisation cannot be assumed. We
repeat the design on seven BoTorch problems with a published maximum
hypervolume --- BraninCurrin, ZDT1, DTLZ2, DH2, VehicleSafety, Penicillin and
CarSideImpact --- spanning two to four objectives and two to seven dimensions,
with the four hypervolume acquisitions (qEHVI, qNEHVI and their log variants)
and the same two model-free floors, at five seeds: $1{,}680$ paired runs, of
which $960$ are model-based noisy-versus-clean comparisons on the problems that
survive the headroom screen below.

Standardisation transfers exactly rather than approximately. Each objective is
mapped to zero mean and unit variance over the same fixed Sobol sample used for
the scalar suite, and under a per-objective affine map the hypervolume scales by
$\prod_m \sigma_m$, so the standardised maximum follows in closed form from the
published one. We verify this rather than assert it: across the seven problems
the largest relative error between the closed-form standardised maximum and a
direct recomputation is $4.7 \times 10^{-16}$.

\paragraph{Controls.} The model-free floor is again identically zero across all
$560$ such comparisons, so no observation noise leaks into the candidate stream
in the hypervolume code path either. No run beats its published maximum
hypervolume. Acquisition-optimisation fallbacks total zero, which we report
rather than assume: a hypervolume acquisition that fails to construct falls back
to random sampling, and because that produces a plausible-looking curve rather
than an error, the arm has to be read off the fallback counter and not off the
plots.

\paragraph{One problem is dropped, by the same screen as before.} DH2's
reference point leaves a $32{,}768$-point Sobol sample with a hypervolume of
exactly zero, and its best learner does not improve on the model-free floor at
all. There is nothing for error to destroy, so its null is guaranteed by
construction; it is excluded. The remaining six pass, with Penicillin weakest at
$0.197$ --- comparable to Rosenbrock at $0.176$, the weakest admissible scalar
landscape.

\begin{table}[t]
\centering
\caption{The two arms on a common footing. Both columns are post-onset excess
divided by the gap between the optimum and a same-budget model-free floor, as a
percentage; brackets are 95\% cluster-bootstrap intervals over
landscapes/problems. This is \emph{not} the normalisation used in
Table~\ref{tab:dose}, and the two must not be read against each other; the
scalar column pools all four error processes, so it also differs from the
\textsc{gaussian}-only column of Table~\ref{tab:fitted}.}
\label{tab:mo}
\begin{tabular}{lrrrr}
\toprule
& \multicolumn{2}{c}{scalar (20 landscapes)} & \multicolumn{2}{c}{multi-objective (6 problems)} \\
\cmidrule(lr){2-3} \cmidrule(lr){4-5}
error & \multicolumn{1}{c}{from trial 1} & \multicolumn{1}{c}{from trial 21} & \multicolumn{1}{c}{from trial 1} & \multicolumn{1}{c}{from trial 21} \\
\midrule
$0.05\sigma$ & 10.2 \makebox[3.4em][l]{\scriptsize [4, 19]} & 2.8 \makebox[3.4em][l]{\scriptsize [1, 5]} & 1.4 \makebox[3.4em][l]{\scriptsize [0, 3]} & 0.5 \makebox[3.4em][l]{\scriptsize [$-$1, 2]} \\
$0.25\sigma$ & 44.3 \makebox[3.4em][l]{\scriptsize [22, 71]} & 6.5 \makebox[3.4em][l]{\scriptsize [4, 10]} & 13.6 \makebox[3.4em][l]{\scriptsize [8, 19]} & 3.7 \makebox[3.4em][l]{\scriptsize [0, 10]} \\
$1\sigma$ & 99.0 \makebox[3.4em][l]{\scriptsize [54, 154]} & 10.3 \makebox[3.4em][l]{\scriptsize [7, 13]} & 49.1 \makebox[3.4em][l]{\scriptsize [42, 56]} & 10.7 \makebox[3.4em][l]{\scriptsize [2, 24]} \\
$5\sigma$ & 148.4 \makebox[3.4em][l]{\scriptsize [76, 239]} & 16.5 \makebox[3.4em][l]{\scriptsize [11, 22]} & 72.3 \makebox[3.4em][l]{\scriptsize [55, 91]} & 13.0 \makebox[3.4em][l]{\scriptsize [3, 27]} \\
\bottomrule
\end{tabular}

\end{table}

\paragraph{The comparison needs one piece of care.} Elsewhere the paper divides
excess regret by $\optz$, a z-score of the optimum against a random design's
mean. A Pareto problem has no scalar optimum and therefore no $\optz$, and the
natural hypervolume analogue --- the gap between the published maximum and what a
model-free design of the same budget reaches --- is a numerically much smaller
denominator. Quoting one arm's published fraction against the other's would
manufacture a threefold difference out of the choice of divisor alone. Both
columns of Table~\ref{tab:mo} are therefore recomputed on the floor denominator,
identically; and we also report the onset ratio, which is a ratio of two costs on
the same problem and so does not depend on the denominator at all.

\textbf{The two headline structures replicate.} Cost rises monotonically with
error magnitude and falls sharply when the error arrives late, on both arms. At
$1\sigma$ from the first rating the multi-objective arm loses $49.1\%$ of the hypervolume floor gap $[42, 56]$ against the scalar arm's $93.7\%$
$[52, 145]$; arriving at iteration~21 the two are indistinguishable, $10.7\%$
$[2, 24]$ against $10.7\%$ $[7, 14]$.

\textbf{There is no evidence multi-objective BO is the more fragile of the two,
and some that it is less.} The point estimates for early error are roughly half
the scalar arm's at every magnitude, and the intervals separate at $0.05\sigma$
and $0.25\sigma$; at $1\sigma$ they overlap slightly and at $5\sigma$
substantially, so this is a consistent direction rather than a
clean separation at the top of the range. The onset ratio is correspondingly
flatter: $3.0$--$5.5\times$ across magnitudes against the scalar arm's
$5.0$--$8.9\times$. A mechanism consistent with both observations, and with
Appendix~\ref{sec:incumbent}, is that a hypervolume objective has no single
incumbent to corrupt --- the state a noisy observation can damage is a Pareto
set, and an error large enough to promote one dominated point does not discard
the rest of the front.

\begin{table}[t]
\centering
\caption{Hypervolume acquisitions, ranked over problems within each condition.
Absolute loss is the fraction of the hypervolume floor gap lost under
noise. The four are not distinguishable.}
\label{tab:moacq}
\begin{tabular}{lrr}
\toprule
acquisition & mean rank & absolute loss \\
\midrule
qNEHVI & 2.40 & 20.2\% \\
qLogEHVI & 2.44 & 20.4\% \\
qEHVI & 2.56 & 20.7\% \\
qLogNEHVI & 2.60 & 20.8\% \\
\bottomrule
\end{tabular}

\end{table}

\textbf{The choice among hypervolume acquisitions does not appear to matter.}
Mean ranks span $2.40$ to $2.60$ on a scale where $2.5$ is a tie, absolute losses
span $20.2\%$ to $20.8\%$, and the Friedman test reaches $p < 0.05$ in $0$ of
$8$ conditions (Table~\ref{tab:moacq}). Kendall's $W$ over
condition-averaged losses is $0.056$, against $0.52$ for the scalar suite computed
the same way, and each of the four acquisitions is the most
robust on at least one of the six problems. Six problems at five seeds is far
less power than twenty at ten, so this is ``no detectable
difference'' rather than ``no difference''. It does, though, point the same way as
Section~\ref{sec:acq}: the disagreement between design spaces about which
acquisition is most robust is, if anything, sharper here.

\paragraph{What does not transfer.} Equation~\ref{eq:frag} is defined for a
scalar objective: it is the loss from one greedy pick among $m$ candidates, and
on a Pareto problem there is no single greedy pick whose loss it measures. The
mediation result is therefore untested in this arm rather than confirmed or
refuted by it. A hypervolume analogue would have to define the one-shot loss over
hypervolume contributions and take a stand on how per-objective error propagates
into them; that is a paper of its own, and we do not claim
Table~\ref{tab:mediation} extends until someone writes it.

\FloatBarrier
\section{Three design choices that could have produced the result}
\label{sec:designchoices}

Every number above was measured at a fixed budget, against an unbounded
continuous response, using standard acquisition functions. Each of those is a
choice, and each could be doing the work.

\begin{table}[t]
\centering
\caption{Budget, with the late onset held at $40\%$ of the run in each case so
that budget is the only thing that changes. Percentages are the fraction of the
achievable improvement destroyed.}
\label{tab:budget}
\begin{tabular}{llrrrr}
\toprule
budget & error present & $0.05\sigma$ & $0.25\sigma$ & $1\sigma$ & $5\sigma$ \\
\midrule
$T=25$ & from it.\ 1 & 0.6 & 3.2 & 9.4 & 15.4 \\
 & from it.\ 11 & 0.3 & 1.3 & 4.1 & 8.0 \\
 & ratio & 2.2$\times$ & 2.4$\times$ & 2.3$\times$ & 1.9$\times$ \\
\midrule
$T=50$ & from it.\ 1 & 0.6 & 3.8 & 11.7 & 21.3 \\
 & from it.\ 21 & 0.1 & 1.0 & 2.3 & 4.7 \\
 & ratio & 4.4$\times$ & 3.9$\times$ & 5.1$\times$ & 4.6$\times$ \\
\midrule
$T=100$ & from it.\ 1 & 0.2 & 3.4 & 11.1 & 23.9 \\
 & from it.\ 41 & 0.1 & 0.4 & 0.9 & 2.2 \\
 & ratio & 1.4$\times$ & 9.2$\times$ & 11.7$\times$ & 10.8$\times$ \\
\bottomrule
\end{tabular}

\end{table}

\paragraph{Budget, and the onset effect is not budget-free.} Re-running at $T=25$
and $T=100$ with the onset held at the same fraction of the run gives
Table~\ref{tab:budget}. Early-error cost is roughly flat in budget --- $9.4\%$,
$14.1\%$, $11.1\%$ at $1\sigma$ --- but late-error cost falls steeply, from
$4.1\%$ at $T=25$ to $2.8\%$ at $T=50$ to $0.9\%$ at $T=100$. The onset ratio
therefore grows with the budget, from $2.3\times$ to $5.0\times$ to $11.7\times$.
This is the clearest limitation on the headline: \textbf{the $5\times$ onset
effect is a statement about a fifty-iteration study.} The direction is
unsurprising, since more remaining iterations means more chance to recover, but
the size is not: a fourfold change in budget moves the ratio by five. We looked
for a rescaling that would let one budget speak for another and did not find
one. Late-onset cost per remaining iteration falls by $18\times$ across the three
budgets, and late-onset cost times remaining iterations is non-monotone
($0.62$, $0.84$, $0.54$), so neither is the invariant. The early-error number is
the transferable one; a study at a different budget has to measure its own onset
effect.

\begin{table}[t]
\centering
\caption{A bounded, discrete response. The instrument arm clips to $\pm 8\sigma$
and rounds to $0.55\sigma$, the range and step implied by the three archival
studies.}
\label{tab:instrument}
\begin{tabular}{llrrrr}
\toprule
response scale & error present & $0.05\sigma$ & $0.25\sigma$ & $1\sigma$ & $5\sigma$ \\
\midrule
unbounded, continuous & from it.\ 1 & 0.6 & 3.8 & 11.7 & 21.3 \\
 & from it.\ 21 & 0.1 & 1.0 & 2.3 & 4.7 \\
bounded, discrete & from it.\ 1 & 3.8 & 5.0 & 11.7 & 20.8 \\
 & from it.\ 21 & 1.3 & 1.4 & 2.5 & 4.6 \\
\bottomrule
\end{tabular}

\end{table}

\paragraph{A real rating scale, which cuts both ways.} The synthetic objective is
unbounded and continuous; a person returns a bounded, discrete rating. The three
archival instruments span $13$, $17$ and $22\sigma$ with steps of $0.55$, $0.54$
and $1.12\sigma$, so the arm clips to $\pm 8\sigma$ and rounds to $0.55\sigma$ ---
the real instrument, in the study's own currency. Table~\ref{tab:instrument}
shows the two effects separately. At $0.05\sigma$ the discretised arm is four
times worse ($3.8\%$ against $0.9\%$), because rounding to a $0.55\sigma$ grid is
itself an error of about $0.16\sigma$ and swamps an injected error far smaller
than the step. At $5\sigma$ it is slightly better ($20.8\%$ against $22.0\%$),
because clipping bounds how far one bad rating can move the incumbent. A rating
scale is a noise floor at the bottom of the range and a safety rail at the top.

\begin{table}[t]
\centering
\caption{Robust baselines. qKG is the knowledge gradient; replication spends
every second evaluation re-asking about the incumbent. Compared at a fixed number
of evaluations, so replication buys averaging with half the designs.}
\label{tab:robust}
\begin{tabular}{llrrrr}
\toprule
acquisitions & error present & $0.05\sigma$ & $0.25\sigma$ & $1\sigma$ & $5\sigma$ \\
\midrule
ten standard & from it.\ 1 & 1.1 & 5.6 & 13.9 & 21.5 \\
 & from it.\ 21 & 0.3 & 1.3 & 2.7 & 4.5 \\
qKG and replication & from it.\ 1 & 0.3 & 2.2 & 9.2 & 16.2 \\
 & from it.\ 21 & 0.2 & 0.1 & 2.2 & 4.8 \\
\bottomrule
\end{tabular}

\end{table}

\paragraph{Methods built for the problem recover about a third.} All twelve
acquisitions above are ordinary ones run under noise, which says which standard
choice survives best but not how much a purpose-built method would recover. Two
baselines: the knowledge gradient, which values the improvement in the posterior
\emph{maximum} rather than in any observed value, and the practical answer of
spending half the budget re-asking about the incumbent. At $1\sigma$ from the
first rating they lose $9.2\%$ of the achievable gain against the standard
acquisitions' $14.1\%$, and at $5\sigma$, $16.2\%$ against $22.0\%$
(Table~\ref{tab:robust}). At the late onset the two are close to
indistinguishable. So a third of the early-error cost is recoverable by choosing
a different method, at a fixed evaluation budget --- which for replication means
buying that recovery with half as many distinct designs. Two thirds is not
recoverable this way. Equation~\ref{eq:frag} measures that residue: the
information the error destroys, rather than any optimizer's failure to cope with
it.

\FloatBarrier
\section{Held-out seeds: a void replication and its sensitivity analysis}
\label{sec:confirmatory}

Everything above was selected and estimated on the same ten seeds. Before
running any fresh ones we wrote down four claims, the statistics that would test
them, the thresholds each had to clear, and four control conditions that would void the run, then ran ten fresh seeds.

\begin{table}[t]
\centering
\caption{Preregistered quantities on the screening seeds and on ten fresh ones.
The replication line reports the \emph{sensitivity} analysis described in the
text, not the preregistered run, which is void.}
\label{tab:confirmatory}
\begin{tabular}{lrr}
\toprule
preregistered quantity & screening & fresh seeds \\
\midrule
H1 $\beta(\frag)$ & +0.65 & +0.69 \\
H1 $\beta(\log_{10}\sigma_e)$ & $-$0.06 & $-$0.07 \\
H2 onset ratio & 5.44 & 5.01 \\
H3 robust mean rank & 4.24 & 4.11 \\
H3 fragile mean rank & 7.26 & 6.98 \\
H4 dose-response monotone & yes & yes \\
\midrule
\multicolumn{3}{l}{replicates: H1: yes; H2: yes; H3: yes; H4: yes} \\
\bottomrule
\end{tabular}

\end{table}

\textbf{The preregistered run is void, by its own rule.} One control failed:
Rosenbrock's headroom on the fresh seeds is $0.083$, below the $0.10$ the
preregistration fixed in advance. The document says that a control failure makes
the run void rather than negative, and an exclusion applied after seeing the data
cannot undo that. We report it as void.

We also report what it showed, labelled as the post-hoc sensitivity analysis it
is. Dropping Rosenbrock from both arms, all four claims clear their thresholds:
$\beta(\frag)$ moves from $+0.65$ to $+0.69$ while $\beta(\log_{10}\sigma_e)$ moves from $-0.06$ to $-0.07$ with an interval
containing zero; the onset ratio is $5.01$ against
$5.44$; the robust and fragile acquisition families sit at mean ranks $4.11$ and
$6.98$ against $4.24$ and $7.26$; and the dose-response is monotone in both
(Table~\ref{tab:confirmatory}).

One argument is worth making explicitly, because it decides how much the void
should worry a reader. The headroom control exists to stop a benchmark with
nothing to lose from contributing a null that looks like evidence of robustness.
It guards against false \emph{negatives}. All four claims came out positive, and
including a low-headroom benchmark makes a positive result harder to obtain, not
easier. The failure mode the control protects against is therefore not one that
can explain what we observed. That is a reason to read the sensitivity analysis
as informative; it is not a reason to call the run confirmed, and we do not.

\FloatBarrier
\section{The landscapes and the descriptor fit}
\label{app:descriptors}

Table~\ref{tab:benchmarks} lists the twenty landscapes with their measured
geometry. Section~\ref{sec:mediation} reports the descriptor model only as a
comparison for $\frag(\sigma_e)$; Table~\ref{tab:descriptors} gives it in full,
with the variance inflation factors that limit what can be read off it.

\begin{table}[htbp]
\centering
\small
\caption{The twenty landscapes and their measured geometry. $\mathrm{opt}_z$ is
how many landscape standard deviations the optimum stands above a random design;
sparsity is the fraction of the box within 10\% of the optimum; ruggedness is
$1$ minus the lag-1 autocorrelation along random walks; tail is
$\sigma/(1.4826\,\mathrm{MAD})$; $\mathrm{frag}(1)$ is the one-shot selection
loss at an error of one landscape standard deviation. Names are the suite's
keys; \texttt{levy\_10} has eleven inputs.}
\label{tab:benchmarks}
\begin{tabular}{lrrrrrrr}
\toprule
landscape & $d$ & $\mathrm{opt}_z$ & sparsity & rugged. & skew & tail & $\mathrm{frag}(1)$ \\
\midrule
\texttt{powell} & 4 & 0.76 & 7.2e-01 & 0.13 & $-$2.33 & 2.52 & 0.18 \\
\texttt{rosenbrock} & 4 & 1.03 & 4.9e-01 & 0.13 & $-$1.06 & 1.13 & 0.18 \\
\texttt{branin} & 2 & 1.25 & 2.4e-01 & 0.13 & $-$1.09 & 1.14 & 0.30 \\
\texttt{steering\_law} & 4 & 1.29 & 7.9e-01 & 0.12 & $-$3.33 & 1.44 & 0.43 \\
\texttt{power\_law\_practice} & 4 & 1.72 & 1.4e-01 & 0.13 & $-$2.02 & 1.72 & 0.63 \\
\texttt{weber\_fechner} & 4 & 2.15 & 5.5e-02 & 0.15 & $-$0.71 & 1.01 & 0.64 \\
\texttt{levy\_10} & 11 & 2.62 & 1.3e-02 & 0.20 & $-$0.54 & 1.00 & 0.80 \\
\texttt{griewank} & 6 & 2.75 & 1.0e-02 & 0.12 & $-$0.26 & 0.98 & 0.85 \\
\texttt{moving\_peaks} & 4 & 2.96 & 1.8e-03 & 0.17 & $-$0.53 & 1.04 & 0.64 \\
\texttt{hartmann\_3} & 3 & 3.05 & 1.6e-02 & 0.14 & +1.14 & 1.49 & 0.50 \\
\texttt{stevens} & 4 & 3.06 & 3.0e-03 & 0.12 & $-$0.17 & 0.97 & 0.75 \\
\texttt{eggholder} & 2 & 3.21 & 6.6e-03 & 0.42 & +0.04 & 1.00 & 0.68 \\
\texttt{yerkes\_dodson} & 4 & 3.54 & 1.9e-03 & 0.12 & +0.27 & 0.98 & 0.73 \\
\texttt{rastrigin} & 4 & 3.64 & 5.8e-04 & 0.58 & $-$0.08 & 0.98 & 0.69 \\
\texttt{schwefel} & 4 & 4.33 & 1.2e-04 & 0.25 & +0.00 & 0.99 & 0.69 \\
\texttt{hicks\_law} & 4 & 5.40 & $<$1.5e-05 & 0.12 & +0.40 & 0.97 & 1.00 \\
\texttt{hartmann\_6} & 6 & 7.97 & 9.2e-05 & 0.14 & +2.67 & 2.74 & 0.53 \\
\texttt{michalewicz} & 5 & 8.05 & $<$1.5e-05 & 0.36 & +1.02 & 0.92 & 0.80 \\
\texttt{ackley} & 4 & 20.10 & $<$1.5e-05 & 0.80 & +3.26 & 1.67 & 0.18 \\
\texttt{shekel} & 4 & 56.77 & $<$1.5e-05 & 0.13 & +3.72 & 1.63 & 0.00 \\
\bottomrule
\end{tabular}

\end{table}


\begin{table}[ht]
\centering
\caption{Post-onset excess regret on landscape descriptors and the error
magnitude, predictors $z$-scored, cluster-bootstrapped over landscapes with
Benjamini--Hochberg correction. VIF is computed on the design matrix.}
\label{tab:descriptors}
\begin{tabular}{lrcrr}
\toprule
term & $\beta$ & 95\% CI & $p_{\mathrm{FDR}}$ & VIF \\
\midrule
$\log\sigma_e$ & +0.37 & [+0.15, +0.66] & 0.001$^{**}$ & 1.0 \\
$\log\mathrm{opt}_z$ & +0.64 & [+0.24, +0.74] & 0.001$^{**}$ & 9.2 \\
$\log$ tail & +0.20 & [+0.03, +0.33] & 0.008$^{**}$ & 1.2 \\
rugged. & $-$0.09 & [$-$0.38, +0.12] & 0.519 & 1.3 \\
$d$ & $-$0.03 & [$-$0.26, +0.12] & 0.768 & 1.0 \\
\midrule
\multicolumn{5}{l}{$R^2 = 0.526$, 640 cells, 20 landscape clusters} \\
\bottomrule
\end{tabular}

\end{table}

$\optz$ carries a VIF of $9.2$ and regresses on the other descriptors with
$R^2 = 0.89$, so its coefficient and those of sparsity and skew (excluded here,
VIF $5.8$ and $6.3$) are not separately identified. This is why
Appendix~\ref{sec:manipulation} builds landscapes that vary these properties
independently rather than trying to disentangle them by regression. Dimension
and ruggedness are null here, and Appendix~\ref{sec:manipulation} finds dimension
nearly inert on excess regret once $\optz$ is held fixed, although it does make
the final design harder to identify.

\FloatBarrier
\section{Input error: the design that ships}
\label{app:inputerror}

Under an unnoticed slip the optimizer's record holds the proposal $x$, not the
point $x'$ it evaluated, and the design a practitioner deploys is the recorded
one. Table~\ref{tab:inputdeployed} scores both at the end of the run: the best
design the person actually experienced, and the design the log recommends,
scored where the log says it is.

\begin{table}[ht]
\centering
\caption{Where a slip's cost comes from, at both onsets. \emph{floor}: the
model-free arms, which evaluate displaced points but learn nothing;
\emph{logged}: a learner whose record holds the point actually evaluated;
\emph{unnoticed}: the same learner with the record holding the proposal;
\emph{mislabelling}: the paired difference of the last two, on shared seeds and
slip draws. The early-onset rows also appear in Table~\ref{tab:inputerror}.}
\label{tab:inputmech}
\begin{tabular}{lrrrr}
\toprule
\textsc{slip} & \multicolumn{1}{c}{floor} & \multicolumn{1}{c}{logged} & \multicolumn{1}{c}{unnoticed} & \multicolumn{1}{c}{mislabelling} \\
\midrule
\multicolumn{5}{l}{\textit{error from it.\ 1}} \\
1\% & 0.3 \makebox[3.4em][l]{\scriptsize [0, 1]} & 0.7 \makebox[3.4em][l]{\scriptsize [0, 2]} & 1.5 \makebox[3.4em][l]{\scriptsize [0, 3]} & 0.8 \makebox[3.4em][l]{\scriptsize [0, 2]} \\
5\% & 0.2 \makebox[3.4em][l]{\scriptsize [$-$1, 1]} & 2.8 \makebox[3.4em][l]{\scriptsize [2, 4]} & 5.4 \makebox[3.4em][l]{\scriptsize [3, 8]} & 2.7 \makebox[3.4em][l]{\scriptsize [1, 4]} \\
15\% & 1.2 \makebox[3.4em][l]{\scriptsize [$-$1, 3]} & 7.1 \makebox[3.4em][l]{\scriptsize [5, 10]} & 11.5 \makebox[3.4em][l]{\scriptsize [8, 15]} & 4.4 \makebox[3.4em][l]{\scriptsize [2, 6]} \\
40\% & 6.4 \makebox[3.4em][l]{\scriptsize [0, 13]} & 19.6 \makebox[3.4em][l]{\scriptsize [12, 28]} & 23.2 \makebox[3.4em][l]{\scriptsize [16, 31]} & 3.6 \makebox[3.4em][l]{\scriptsize [2, 6]} \\
\midrule
\multicolumn{5}{l}{\textit{error from it.\ 21}} \\
1\% & 0.0 \makebox[3.4em][l]{\scriptsize [0, 0]} & 0.2 \makebox[3.4em][l]{\scriptsize [0, 1]} & 1.1 \makebox[3.4em][l]{\scriptsize [0, 2]} & 0.9 \makebox[3.4em][l]{\scriptsize [0, 1]} \\
5\% & $-$0.4 \makebox[3.4em][l]{\scriptsize [$-$1, 0]} & 2.0 \makebox[3.4em][l]{\scriptsize [1, 3]} & 3.0 \makebox[3.4em][l]{\scriptsize [2, 5]} & 1.0 \makebox[3.4em][l]{\scriptsize [0, 2]} \\
15\% & $-$0.4 \makebox[3.4em][l]{\scriptsize [$-$1, 1]} & 4.5 \makebox[3.4em][l]{\scriptsize [2, 7]} & 5.0 \makebox[3.4em][l]{\scriptsize [3, 8]} & 0.5 \makebox[3.4em][l]{\scriptsize [0, 1]} \\
40\% & $-$0.6 \makebox[3.4em][l]{\scriptsize [$-$2, 1]} & 6.3 \makebox[3.4em][l]{\scriptsize [3, 10]} & 6.4 \makebox[3.4em][l]{\scriptsize [4, 10]} & 0.1 \makebox[3.4em][l]{\scriptsize [0, 0]} \\
\bottomrule
\end{tabular}

\end{table}

\begin{table}[ht]
\centering
\caption{Final-iteration excess simple regret under an unnoticed slip, as a
percentage of the achievable improvement, over the model-based acquisitions.
\emph{Evaluated}: the best true value among the points actually touched.
\emph{Deployed}: the true value at the recorded location of the best-rated
trial. Brackets are 95\% cluster-bootstrap intervals over landscapes.}
\label{tab:inputdeployed}
\begin{tabular}{lrrrr}
\toprule
& \multicolumn{2}{c}{from trial 1} & \multicolumn{2}{c}{from trial 21} \\
\cmidrule(lr){2-3} \cmidrule(lr){4-5}
\textsc{slip} & \multicolumn{1}{c}{evaluated} & \multicolumn{1}{c}{deployed} & \multicolumn{1}{c}{evaluated} & \multicolumn{1}{c}{deployed} \\
\midrule
1\% & 1.7 \makebox[3.4em][l]{\scriptsize [0, 3]} & 2.7 \makebox[3.4em][l]{\scriptsize [1, 5]} & 1.5 \makebox[3.4em][l]{\scriptsize [1, 3]} & 2.4 \makebox[3.4em][l]{\scriptsize [1, 4]} \\
5\% & 6.5 \makebox[3.4em][l]{\scriptsize [4, 10]} & 16.4 \makebox[3.4em][l]{\scriptsize [9, 24]} & 3.9 \makebox[3.4em][l]{\scriptsize [2, 6]} & 9.6 \makebox[3.4em][l]{\scriptsize [5, 15]} \\
15\% & 13.4 \makebox[3.4em][l]{\scriptsize [8, 20]} & 34.3 \makebox[3.4em][l]{\scriptsize [24, 45]} & 6.9 \makebox[3.4em][l]{\scriptsize [3, 10]} & 13.4 \makebox[3.4em][l]{\scriptsize [8, 19]} \\
40\% & 24.9 \makebox[3.4em][l]{\scriptsize [16, 34]} & 65.8 \makebox[3.4em][l]{\scriptsize [53, 78]} & 9.0 \makebox[3.4em][l]{\scriptsize [5, 14]} & 13.3 \makebox[3.4em][l]{\scriptsize [8, 19]} \\
\bottomrule
\end{tabular}

\end{table}

\FloatBarrier
\section{Does noise ever help?}
\label{app:beneficial}

A noisy observation can in principle help, by pushing an over-exploiting
optimizer out of a local basin. At one seed per cell, $25$ of the $640$
landscape-by-condition cells had negative mean excess regret with $d_z < -0.2$,
noise apparently improving the outcome, which is the kind of finding that
invites a mechanism.

At ten seeds it is one cell: \texttt{michalewicz} under \textsc{ar}(1) error at
$0.05\sigma$ from iteration $21$, mean excess $-0.136$, $d_z = -0.21$ over $100$
paired runs. One cell in $640$ is fewer than the fifteen or so that a null family of $640$
tests would flag at this threshold. We report it as a null and note the
sampling-variability lesson: the twenty-four cells that disappeared between one
seed and ten were all in the direction that would have made the better story.

\end{document}
